%=============================================================================
\section{Sensor Calibration and Selection}
%=============================================================================

Your filters from the previous chapter can remove noise, but they cannot fix a sensor that is fundamentally lying to you. A gyroscope with a $2\textdegree{}/\mathrm{s}$ bias will drift $120\textdegree{}$ per minute no matter how good your low-pass filter is. This chapter teaches you how to make your sensors honest before feeding their data into the control pipeline.

Calibration is not glamorous work. It is the kind of thing you do on a Saturday afternoon with a flat table and a protractor. But it is the single highest-return activity in robotics: ten minutes of calibration can eliminate errors that no amount of software can fix.

%-----------------------------------------------------------------------------
\subsection{The Sensor Error Model}
%-----------------------------------------------------------------------------

Every sensor you will ever use obeys roughly the same error model. The measurement $y_m$ is related to the true value $y$ by:

\begin{equation}
y_m = S \, y + b + n
\end{equation}

where $S$ is the \textbf{scale factor} (ideally 1, but in practice 0.98--1.02 for a typical MEMS sensor), $b$ is the \textbf{bias} (a constant offset that persists even when the true value is zero), and $n$ is \textbf{noise} (zero-mean random fluctuation).

Let us build intuition for each term. The \textbf{bias} is the most insidious: if your gyroscope reads $0.5\textdegree{}/\mathrm{s}$ when the robot is perfectly still, that half-degree-per-second adds up relentlessly. After one minute you have $30\textdegree{}$ of phantom rotation. After five minutes, $150\textdegree{}$. No filter can remove a DC offset from a rate sensor---you must measure it and subtract it.

The \textbf{scale factor} error is subtler. If $S = 1.03$ instead of $1.0$, then a $90\textdegree{}$ turn reads as $92.7\textdegree{}$. You will not notice this during a quick test, but it accumulates over a full match.

The \textbf{noise} $n$ is what you already know how to handle: low-pass filters, complementary filters, Kalman filters. It is the easiest error to deal with.

For multi-axis sensors (accelerometers, gyroscopes, magnetometers), the model generalizes to:

\begin{equation}
\mathbf{y}_m = \mathbf{S} \, \mathbf{y} + \mathbf{b} + \mathbf{n}
\end{equation}

where $\mathbf{S}$ is now a $3 \times 3$ matrix. The off-diagonal elements of $\mathbf{S}$ capture \textbf{cross-axis coupling}---the $x$-axis accelerometer picking up a bit of $y$-axis acceleration due to imperfect mechanical alignment inside the chip. For cheap MEMS sensors, cross-axis coupling is typically 1--3\% of the main axis, which is small but measurable.

\begin{center}
\begin{tabular}{llll}
\toprule
\textbf{Error source} & \textbf{Typical magnitude} & \textbf{How to detect} & \textbf{How to fix} \\
\midrule
Bias (accel.) & $\pm 50\;\mathrm{mg}$ & Static test: read at rest & Subtract mean \\
Bias (gyro.) & $\pm 1\textdegree{}/\mathrm{s}$ & Static test: read at rest & Subtract mean \\
Scale factor & 1--3\% & Compare to known reference & Solve calibration eqn. \\
Cross-axis & 1--3\% & 6-position calibration & Full $3\times3$ matrix \\
Noise (accel.) & $\sim\!200\;\mu g/\!\sqrt{\mathrm{Hz}}$ & Allan variance or PSD & Low-pass filter \\
Noise (gyro.) & $\sim\!0.01\textdegree{}/\mathrm{s}/\!\sqrt{\mathrm{Hz}}$ & Allan variance or PSD & Low-pass filter \\
Temp. drift & varies & Test at different temps & Online compensation \\
\bottomrule
\end{tabular}
\end{center}

Temperature drift deserves a mention. MEMS sensor biases shift with temperature---sometimes significantly. High-end IMUs include internal temperature sensors and apply polynomial compensation. For competition robots operating in a gym at roughly constant temperature, you can usually ignore this, but be aware that a robot calibrated indoors in winter may behave differently at an outdoor summer event.

%-----------------------------------------------------------------------------
\subsection{IMU Calibration}
%-----------------------------------------------------------------------------

\subsubsection{Accelerometer: 6-Position Static Calibration}

The idea is beautifully simple: gravity gives you a known, free, highly accurate reference vector. On Earth's surface, $\|\mathbf{g}\| = 9.81\;\mathrm{m/s^2}$ to four significant figures. If you can orient your accelerometer so that gravity aligns with each axis in turn (positive and negative), you get six data points that fully constrain the $3 \times 3$ scale-and-cross-axis matrix $\mathbf{S}$ and the $3 \times 1$ bias vector $\mathbf{b}$.

Place your IMU on a flat surface in six orientations: $+x$ up, $-x$ up, $+y$ up, $-y$ up, $+z$ up, $-z$ up. In each orientation, collect a few hundred samples and average them. You do not need a precision jig---a flat table and a box with square edges work fine.

In each orientation $i$, the measurement model gives:

\begin{equation}
\mathbf{a}_{m,i} = \mathbf{S} \, \mathbf{g}_i + \mathbf{b}
\end{equation}

where $\mathbf{g}_i$ is the known gravity vector in that orientation. For example, when $+z$ is up, $\mathbf{g}_i = [0, 0, 9.81]^T$. When $-x$ is up, $\mathbf{g}_i = [-9.81, 0, 0]^T$.

Stack all six equations into a single matrix equation. Define $\mathbf{A}_m \in \mathbb{R}^{6 \times 3}$ as the matrix of measured accelerations (one row per orientation) and $\mathbf{G} \in \mathbb{R}^{6 \times 3}$ as the corresponding known gravity vectors. Then:

\begin{equation}
\mathbf{A}_m = \mathbf{G} \, \mathbf{S}^T + \mathbf{1} \, \mathbf{b}^T
\end{equation}

This is a standard linear system. Augment $\mathbf{G}$ with a column of ones to absorb the bias, and solve via least squares:

\begin{equation}
\begin{bmatrix} \mathbf{S}^T \\ \mathbf{b}^T \end{bmatrix}
= \left( \tilde{\mathbf{G}}^T \tilde{\mathbf{G}} \right)^{-1} \tilde{\mathbf{G}}^T \mathbf{A}_m
\end{equation}

where $\tilde{\mathbf{G}} = [\mathbf{G} \;\; \mathbf{1}]$ is the augmented matrix. Once you have $\mathbf{S}$ and $\mathbf{b}$, the corrected reading is:

\begin{equation}
\mathbf{a}_{\text{true}} = \mathbf{S}^{-1} (\mathbf{a}_m - \mathbf{b})
\end{equation}

Here is a minimal Python script that performs the entire calibration:

\begin{verbatim}
import numpy as np

# a_meas: (6, 3) array of mean accelerometer readings per orientation
# g_ref:  (6, 3) array of known gravity vectors per orientation
G_aug = np.hstack([g_ref, np.ones((6, 1))])          # (6, 4)
params, _, _, _ = np.linalg.lstsq(G_aug, a_meas, rcond=None)
S = params[:3, :].T                                   # (3, 3)
b = params[3, :]                                       # (3,)
S_inv = np.linalg.inv(S)

def calibrate(a_raw):
    return S_inv @ (a_raw - b)
\end{verbatim}

After calibration, verify: rotate the IMU to arbitrary orientations and check that $\|\mathbf{a}_{\text{true}}\| \approx 9.81\;\mathrm{m/s^2}$. If the magnitude wanders by more than $\pm 0.05\;\mathrm{m/s^2}$, something went wrong---most likely one of your six orientations was not level.

\subsubsection{Gyroscope: Static Bias and Noise}

Gyroscope calibration is simpler than accelerometer calibration because you only need one test condition: \textit{be perfectly still}.

Place your IMU on a stable surface and record data for 60 seconds. The average reading on each axis is your bias. Subtract it in your code:

\begin{verbatim}
# During initialization (robot stationary):
gyro_bias = np.mean(gyro_samples, axis=0)    # average over ~60s
# During operation:
gyro_corrected = gyro_raw - gyro_bias
\end{verbatim}

That is the entire calibration for most competition applications. Do this every time you power on, because MEMS gyroscope bias drifts slightly from session to session (``turn-on bias'').

For a deeper characterization, the \textbf{Allan variance} tells you more about your sensor's noise structure. The Allan variance is computed by averaging the gyroscope output over progressively longer time windows $\tau$ and plotting the resulting standard deviation $\sigma(\tau)$ on a log-log plot.

The key features of the Allan variance plot are:

At short averaging times ($\tau < 1\;\mathrm{s}$), the curve slopes downward with slope $-1/2$. This region is dominated by \textbf{angle random walk} (white noise)---the fundamental noise floor of the sensor. The value of $\sigma$ at $\tau = 1\;\mathrm{s}$ gives you the angle random walk coefficient directly.

At intermediate averaging times, the curve reaches a minimum. This minimum is the \textbf{bias instability}---the best stability your sensor can achieve. For a \$2 MEMS gyroscope, this is typically around $1$--$10\textdegree{}/\mathrm{hr}$; for a fiber-optic gyroscope, it might be $0.001\textdegree{}/\mathrm{hr}$.

At long averaging times ($\tau > 100\;\mathrm{s}$), the curve turns upward with slope $+1/2$. This is \textbf{rate random walk}---slow drift processes like temperature changes.

For competition robots, do not overthink this. Measure the static bias, subtract it, and move on. Allan variance analysis is useful when you are selecting between sensors or debugging long-term drift, but your matches are 3--5 minutes long, and a simple bias subtraction handles that timescale just fine.

\subsubsection{Gyroscope Scale Factor}

If you need to calibrate the gyroscope scale factor (not just bias), you need a known rotation rate. The poor-man's approach: mount the IMU on a turntable (a lazy Susan works), rotate it exactly $360\textdegree{}$, and check whether the integrated gyroscope output agrees. If the gyro says you rotated $354\textdegree{}$, your scale factor is $360/354 \approx 1.017$. This is rarely necessary for competition robots---MEMS gyroscope scale factors are typically within 1\% of nominal---but it takes five minutes and eliminates one more source of error.

A more precise method uses a rate table that spins at a known, constant angular velocity. You compare the gyroscope reading to the known rate at several speeds, fit a line, and extract the slope (scale factor) and intercept (bias) simultaneously. This is overkill for most competition scenarios, but it is how the sensor manufacturers characterize their parts.

%-----------------------------------------------------------------------------
\subsection{Encoder: Resolution vs Noise Tradeoff}
%-----------------------------------------------------------------------------

\subsubsection{Velocity Estimation from Position}

Encoders give you position (angle), but your control loop almost always needs velocity. The naive approach is a finite difference:

\begin{equation}
v[n] = \frac{\theta[n] - \theta[n-1]}{T_s}
\end{equation}

where $T_s$ is the sample period. This works well at high speeds but falls apart at low speeds due to \textbf{quantization}. An encoder with $N$ counts per revolution has a minimum detectable position change of $\Delta\theta = 2\pi / N$ radians. The minimum nonzero velocity you can measure is therefore:

\begin{equation}
v_{\min} = \frac{2\pi / N}{T_s}
\end{equation}

Let us plug in numbers. A 1024-count encoder sampled at $1\;\mathrm{kHz}$ gives $v_{\min} = 2\pi / (1024 \times 0.001) \approx 6.14\;\mathrm{rad/s} \approx 58.6\;\mathrm{RPM}$. Below this speed, the velocity estimate jumps between 0 and $v_{\min}$ with nothing in between. Your motor appears to be stuttering even though it is spinning smoothly.

There are several solutions, in order of increasing sophistication.

\textbf{Reduce the sample rate for velocity.} If you compute velocity every $10\;\mathrm{ms}$ instead of every $1\;\mathrm{ms}$, $v_{\min}$ drops by a factor of 10. The downside is a 10$\times$ slower velocity update, which may be too slow for your inner control loop.

\textbf{M/T method.} Instead of counting edges per fixed time interval (M method), measure the time between edges (T method), or combine both (M/T method). At low speeds, where edges arrive slowly, measuring the time between them gives much finer velocity resolution. This is what hardware timer capture/compare peripherals are designed for. On an STM32, you configure a timer channel in ``input capture'' mode: every time an encoder edge arrives, the hardware latches the timer counter value. The time between two captures, divided into the angle per edge, gives you velocity with resolution limited only by the timer clock (typically $72$--$168\;\mathrm{MHz}$). The M/T hybrid automatically switches between edge-counting at high speed and time-measurement at low speed, giving you the best of both worlds across the entire speed range.

\textbf{Low-pass filter the finite difference.} A simple first-order LPF on the velocity estimate smooths out the quantization jumps. This introduces a small delay but is often the best tradeoff for simplicity:

\begin{verbatim}
v_raw = (theta - theta_prev) / Ts;
v_filtered += alpha * (v_raw - v_filtered);
theta_prev = theta;
\end{verbatim}

\textbf{Observer-based estimation.} A Luenberger observer or a Kalman filter can estimate velocity from noisy position measurements with better performance than a simple filter. The observer uses a model of the motor dynamics (typically a first-order model: $\dot{\omega} = (K_t i - B\omega) / J$) together with the position measurement to produce a smooth velocity estimate. We will develop this idea properly in Chapter 6 when we discuss state observers, so for now just know that this option exists and is the theoretically optimal approach.

\subsubsection{Encoder Types and Selection}

Not all encoders are created equal. \textbf{Incremental encoders} output two quadrature pulse trains (A and B channels) and optionally an index pulse (Z channel) that fires once per revolution. The quadrature encoding gives you direction as well as position, and most microcontroller timer peripherals can decode quadrature signals in hardware with zero CPU overhead. Incremental encoders lose their position when power is removed, so you need a homing procedure at startup.

\textbf{Absolute encoders} report the exact shaft angle at all times, even after a power cycle. They use either optical or magnetic sensing with a digital output (SPI, SSI, or analog). The AS5048A and AS5600 are popular magnetic absolute encoders in the robotics world. Their resolution is typically 12--14 bits ($4096$--$16384$ counts per revolution), which is excellent. The downside is that communication latency (a few microseconds for SPI) can matter in very fast control loops, and some cheap magnetic encoders have nonlinearity of $\pm 1$--$2$ LSB that varies with shaft angle.

For most competition robots, a 1024--4096 count incremental encoder on the motor shaft (before the gearbox) gives the best combination of resolution, speed, and simplicity. If you need absolute position (e.g., for a joint that must know its angle at startup without homing), use a magnetic absolute encoder on the output shaft.

%-----------------------------------------------------------------------------
\subsection{Magnetometer Calibration}
%-----------------------------------------------------------------------------

A magnetometer measures the Earth's magnetic field to determine heading. In a perfect world, if you rotate the sensor $360\textdegree{}$ in the horizontal plane, the $x$-$y$ readings trace out a perfect circle centered at the origin. In reality, you get an offset ellipse, and this is where calibration comes in.

\textbf{Hard iron} distortion comes from permanent magnets or magnetized ferrous materials near the sensor (steel screws, speaker magnets, battery contacts). These produce a constant magnetic field that shifts the circle away from the origin. The fix is to find the center of the shifted circle and subtract it---that is your hard iron offset vector.

\textbf{Soft iron} distortion comes from magnetically permeable materials (soft steel, iron) that distort the Earth's field non-uniformly, stretching the circle into an ellipse. The fix is to fit an ellipse to the data and compute the transformation matrix that maps it back to a circle.

The practical calibration procedure: slowly rotate your robot $360\textdegree{}$ (or better, tumble it in all orientations), record the magnetometer data, fit an ellipse to the $x$-$y$ scatter plot, and extract the offset and scaling. Many open-source tools do this automatically.

A word of caution: magnetometers are \textbf{unreliable indoors}. Steel structural beams, reinforced concrete, electrical wiring, and nearby motors all distort the magnetic field in unpredictable, location-dependent ways. If your competition is in a building with a steel frame, your magnetometer heading may vary by $10$--$30\textdegree{}$ depending on where the robot is on the field. Use the magnetometer for coarse heading initialization, but do not trust it for continuous navigation indoors. Gyroscope integration (with periodic corrections) is almost always more reliable for short-duration matches.

%-----------------------------------------------------------------------------
\subsection{Current Sensor Calibration}
%-----------------------------------------------------------------------------

If you are doing field-oriented control (FOC) of a BLDC motor, your current measurements feed directly into the innermost control loop. Errors here propagate everywhere.

\textbf{Offset calibration:} Before the motor starts, with zero current flowing, read the current sensor for a few hundred samples and average. This is your zero-current offset. Subtract it from all subsequent readings. On many motor driver boards, the current sensor is a shunt resistor plus an op-amp with a mid-supply reference voltage, so the ``zero'' reading is around $V_{cc}/2$, not actually $0\;\mathrm{V}$.

\textbf{Bandwidth matters.} Your current control loop typically runs at $10$--$20\;\mathrm{kHz}$. If your current sensor has a bandwidth of only $5\;\mathrm{kHz}$, the sensor itself acts as a low-pass filter with significant phase lag at the frequencies where your controller is trying to act. This phase lag reduces your phase margin and can destabilize the current loop. Check your sensor datasheet: the bandwidth should be at least $5\times$ higher than your current loop frequency. Hall-effect current sensors (like the ACS712) are typically limited to a few hundred kHz, which is fine. Shunt-resistor-based sensing with a fast op-amp can reach MHz bandwidth, which is more than enough.

%-----------------------------------------------------------------------------
\subsection{Sensor Fusion Architecture}
%-----------------------------------------------------------------------------

Once your sensors are calibrated, the next question is how to combine them. There are two fundamentally different approaches.

\textbf{Measurement-level fusion} blends the raw (calibrated) sensor data directly. The complementary filter from Chapter 2 is the classic example: low-pass the accelerometer, high-pass the gyroscope, add them together. This is simple, lightweight, and works beautifully when you have exactly two sensors with complementary noise characteristics.

\textbf{State-level fusion} estimates the system's internal state (position, velocity, orientation) from all sensor measurements simultaneously, using a dynamic model of the system. The Kalman filter is the gold standard here. It handles any number of sensors, accounts for their individual noise characteristics, and produces statistically optimal estimates under certain assumptions.

How do you choose? Here is a simple decision rule:

If you have 1--2 sensors with well-understood, complementary noise profiles, and the relationship between sensors and the state of interest is simple, use a complementary filter. It is five lines of code, runs in microseconds, and is easy to tune.

If you have 3 or more sensors, or the relationship between measurements and states is nonlinear (e.g., GPS + IMU + magnetometer for full pose estimation), or you want provably optimal performance, use the Extended Kalman Filter from Chapter 6. The EKF handles all of these cases in a unified framework. The cost is more code, more tuning, and more computation---but modern microcontrollers handle it easily.

Most competition robots land in the complementary filter camp. An accelerometer-gyroscope complementary filter for tilt estimation, perhaps augmented with encoder odometry for position. Save the EKF for when you genuinely need it.

A useful mental model: think of sensor fusion as \textit{spectral division of labor}. The accelerometer is trustworthy at low frequencies (it sees gravity, which is constant) but contaminated by vibration at high frequencies. The gyroscope is trustworthy at high frequencies (it responds instantly to rotation) but drifts at low frequencies (bias accumulation). A complementary filter simply assigns each sensor to the frequency range where it excels. The Kalman filter does the same thing, but it computes the optimal crossover frequency automatically from the noise statistics.

\subsection{Choosing the Right Sensor}

Before you calibrate, you need to pick the right sensor in the first place. Here is a compact decision guide for the most common competition-robot sensing needs.

\textbf{Orientation (tilt/heading):} A 6-axis IMU (accelerometer + gyroscope) covers most needs. Add a magnetometer only if you need absolute heading and your arena does not have large steel structures nearby. The BMI088, ICM-42688, and MPU-6500 are popular choices in order of increasing quality.

\textbf{Wheel velocity:} Use a quadrature encoder with at least $N = 4096$ counts per revolution if you need smooth velocity down to $\sim\!5\;\mathrm{RPM}$ at $1\;\mathrm{kHz}$ sampling. If your motor has a gearbox, count on the motor shaft (before the gearbox) for higher resolution.

\textbf{Distance/ranging:} Time-of-flight (ToF) sensors like the VL53L1X give millimeter resolution up to $\sim\!4\;\mathrm{m}$. They are small, cheap, and easy to use over I2C. For longer range, a 2D lidar (e.g., RPLIDAR A1) gives $360\textdegree{}$ scans at $5$--$10\;\mathrm{Hz}$ up to $12\;\mathrm{m}$.

\textbf{Current sensing for FOC:} Use inline shunt resistors with a fast op-amp for best bandwidth. Hall-effect sensors (ACS712, ACS714) are easier to integrate but have lower bandwidth and higher noise. At a minimum, your current sensor bandwidth should exceed $5\times$ the current loop frequency.

The general principle: buy the best sensor you can afford, then calibrate it carefully. A well-calibrated \$3 sensor will outperform a poorly calibrated \$30 sensor every time.

%-----------------------------------------------------------------------------
\subsection{Practical: Verifying Calibration Quality}
%-----------------------------------------------------------------------------

Calibration is not a ``set and forget'' operation. You need to verify that it actually worked, and you need to log the right data so that you can debug problems later.

\textbf{Accelerometer sanity check:} After calibration, hold the IMU in several arbitrary orientations. Compute $\|\mathbf{a}_{\text{cal}}\|$. It should be $9.81 \pm 0.05\;\mathrm{m/s^2}$ in every orientation. If it deviates, your calibration matrix is wrong---probably because one of the six positions was not level during the calibration procedure.

\textbf{Gyroscope sanity check:} With the robot stationary, integrate the calibrated gyroscope rate over 60 seconds. The accumulated angle should stay within $\pm 1\textdegree{}$. If it drifts further, your bias estimate is off. Re-run the static calibration, making sure the robot is truly motionless (no fans causing vibration, no one bumping the table).

\textbf{Magnetometer sanity check:} After hard/soft iron calibration, rotate the sensor $360\textdegree{}$ in the horizontal plane and plot the calibrated $x$-$y$ data. It should form a circle centered near the origin, not an ellipse or an offset blob.

\textbf{What to log:} Always log both raw and calibrated data with timestamps. When something goes wrong on the field, you need the raw data to determine whether the problem is in the sensor or in the calibration. If you only log calibrated data, you cannot tell the difference. Logging is cheap---a few kilobytes per second. The debugging time it saves is enormous.

\textbf{Encoder sanity check:} Spin the wheel exactly one full revolution by hand (use a mark on the wheel). The encoder should report exactly $N$ counts (or $4N$ in quadrature mode, depending on your counting configuration). If it consistently reads $N-1$ or $N+1$, you may have electrical noise causing spurious edges---add a small capacitor across the encoder output lines, or reduce the cable length.

\textbf{Current sensor sanity check:} With the motor disconnected, the current reading should be zero (after offset subtraction). Then command a known, steady current and verify with a multimeter or a clamp meter. If the sensor reads $1.0\;\mathrm{A}$ but the multimeter says $0.95\;\mathrm{A}$, your gain calibration is off by 5\%.

\textbf{When to recalibrate:} Recalibrate the gyroscope bias every time you power on---it shifts by a few tenths of a degree per second between sessions. Accelerometer and magnetometer calibrations are stable unless you change the mounting or add new ferrous components to the chassis. Current sensor offsets should be measured at every startup, before the motor is energized.

%-----------------------------------------------------------------------------
% Case Studies
%-----------------------------------------------------------------------------

\vspace{0.5em}
\begin{mdframed}[linewidth=1pt, innertopmargin=10pt, innerbottommargin=10pt]
\textbf{Case Study: My Heading Drifts $30\textdegree{}$ in 5 Minutes}

\vspace{0.3em}
The problem: a competition robot uses a magnetometer for heading estimation. During testing, the heading slowly drifts $30\textdegree{}$ over a 5-minute run, despite the robot returning to its starting orientation.

The diagnosis starts with plotting the raw magnetometer $x$-$y$ data while rotating the robot $360\textdegree{}$. Instead of a circle centered at the origin, the data forms a circle centered at roughly $(150, -80)\;\mu\mathrm{T}$. This is classic hard iron distortion---the steel chassis and the neodymium magnets in the drive motors are producing a constant magnetic offset.

The fix: rotate the robot through a full $360\textdegree{}$ turn (slowly, on a flat surface), record the magnetometer data, compute the center of the circle $(c_x, c_y)$, and subtract it from all future readings:

\begin{verbatim}
mag_cal_x = mag_raw_x - cx;
mag_cal_y = mag_raw_y - cy;
heading = atan2(mag_cal_y, mag_cal_x);
\end{verbatim}

After this three-line fix, the heading error over 5 minutes drops from $30\textdegree{}$ to less than $3\textdegree{}$. The remaining error is from soft iron distortion and indoor magnetic field non-uniformity---acceptable for a competition robot that also has gyroscope integration as a backup.

The lesson: always visualize your magnetometer data as a scatter plot before trusting it. If the circle is not centered at the origin, you have a hard iron problem. If it is an ellipse instead of a circle, you also have a soft iron problem.
\end{mdframed}

\vspace{0.5em}
\begin{mdframed}[linewidth=1pt, innertopmargin=10pt, innerbottommargin=10pt]
\textbf{Case Study: Velocity Goes Crazy Below 10 RPM}

\vspace{0.3em}
The problem: a wheeled robot uses a 1024-count quadrature encoder for wheel velocity feedback. The PID velocity controller works well above $60\;\mathrm{RPM}$, but below about $10\;\mathrm{RPM}$ the motor audibly stutters and the velocity plot looks like a square wave oscillating between $0$ and some minimum value.

The analysis: at $1\;\mathrm{kHz}$ sampling, the minimum detectable velocity for a 1024-count encoder is:

\[
v_{\min} = \frac{2\pi}{1024 \times 0.001\;\mathrm{s}} \approx 6.14\;\mathrm{rad/s} \approx 58.6\;\mathrm{RPM}
\]

Below this speed, the position difference $\theta[n] - \theta[n-1]$ is almost always 0 (no edge counted), with occasional jumps of $\pm 1$ count. The finite-difference velocity oscillates between 0 and $v_{\min}$. The PID controller sees ``zero velocity'' and cranks up the output; then it sees ``$v_{\min}$'' and backs off. This creates the audible stuttering.

The fix combines two changes. First, compute velocity at a lower rate---every $10\;\mathrm{ms}$ instead of every $1\;\mathrm{ms}$---which reduces $v_{\min}$ to $5.86\;\mathrm{RPM}$. Second, apply a first-order low-pass filter ($\alpha = 0.1$) to smooth the remaining quantization steps:

\begin{verbatim}
// Every 10 ms:
v_raw = (encoder_count - prev_count) * (2*PI/1024) / 0.01;
v_filt += 0.1 * (v_raw - v_filt);
prev_count = encoder_count;
\end{verbatim}

After this change, the motor spins smoothly down to about $3\;\mathrm{RPM}$. For even lower speeds, the team later switched to the M/T method using the STM32 timer's input capture mode, which measures the time between encoder edges and provides fine velocity resolution all the way down to near-zero speed.

The lesson: always compute $v_{\min}$ for your encoder-and-sample-rate combination \textit{before} building the robot. If $v_{\min}$ is above the minimum speed you need, either use a higher-resolution encoder, reduce the velocity sample rate, or implement M/T measurement.
\end{mdframed}

\vspace{0.5em}
\subsection{Calibration Checklist}

Here is the calibration sequence we recommend running every time you bring up a new robot or make significant hardware changes. It takes about 30 minutes and saves hours of debugging later.

\begin{center}
\begin{tabular}{rll}
\toprule
\textbf{Step} & \textbf{Sensor} & \textbf{Procedure} \\
\midrule
1 & Current sensor & Zero-current offset (motor disconnected) \\
2 & Gyroscope & 60-second static bias (robot stationary on flat surface) \\
3 & Accelerometer & 6-position static calibration (if first time or remount) \\
4 & Magnetometer & $360\textdegree{}$ rotation for hard/soft iron (if used) \\
5 & Encoder & One-revolution count check (verify $N$ counts exactly) \\
6 & All sensors & Log raw + calibrated data, run sanity checks \\
\bottomrule
\end{tabular}
\end{center}

Steps 1 and 2 should be done at every power-on. Steps 3--5 only need to be repeated when hardware changes. Step 6 should become a habit---the five minutes you spend verifying calibration will save you from chasing phantom bugs during the competition.

The bottom line of this entire chapter is simple: \textit{calibrate before you filter, filter before you control.} A sensor that tells the truth, even noisily, is infinitely more useful than one that lies quietly. Get the truth first; the rest of this book will teach you what to do with it.
